\section{Discussion and limitations}
A source-valid memory provides both a procedure and evidence about the circumstances in which it worked. The observed agent behavior is consistent with either component becoming salient. Some implementations preserve the source procedure without the target check; others use the stated source assumption to identify what must change. The successful F08 C runs show why a one-direction account of memory as harmful imitation is incomplete. They also explain why adding a boundary reminder can have little marginal benefit when the source text already supports applicability reasoning.

Three evaluation choices are consequential. First, functionality and the focal witness must be reported jointly. A task left at its baseline may avoid the behavior the security test observes, but it has not met the feature request. Second, a witness must observe the contract-relevant return value or effect. Reference controls alone did not prevent the X05/X28 bookkeeping defects exposed by alternative agent implementations. Third, invalidity can depend on the implementation: an agent may recognize the main authority boundary yet fail on its error path. Treating such runs as unrelated outages would misstate what is missing.

These implications support better memory evaluation, but the experiment has a narrow reach. Thirteen deliberately constructed families with two repetitions per cell provide limited precision and are not a probability sample of software work. The tested memory is supplied directly and does not evaluate retrieval, online learning, or the causal effect of context shift. Prompt bytes are balanced among memory conditions, but tokens and information content are not identical. Tool instructions, provider prompts, and budgets differ across harnesses, so the observed cross-configuration differences cannot be assigned to model weights alone.

The security outcomes are bounded by their executed witnesses. X06 rejection is evidence of confidentiality on one mixed fixture, not full retention or general redaction. The X05 correction observes emitted packets and does not establish arbitrary concurrent allocator behavior. The other families were not exhaustively audited for all possible contract violations. Finally, the behavioral analysis is selected and outcome-aware: visible statements can corroborate implementation changes, but they do not reveal all internal computation or establish causal mediation. The study supplies evidence for conditional transfer and for measurement distinctions that future, larger evaluations should preserve.

\section{Conclusion}
Supplying correct procedural memory under changed trust assumptions produced heterogeneous outcomes in this controlled coding study. The original harmful-memory hypothesis was not established as a general effect. Local harmful transfer, successful adaptation, and non-completion all occurred, and a generic applicability reminder helped descriptively in some configurations without becoming a reliable general repair. The clearest lesson is methodological: evaluate what the agent completed, which security property was actually observed, and how the implementation relates to the visible procedure before drawing conclusions about memory transfer.
